\chapter{\texttt{CompactRelation}: 32 bytes por arista}
\label{cap:compact-relation}

\begin{intuicion}
La \texttt{Relation} pública pesa \(\sim 200\)~bytes en RAM por la
metadata flexible y los IDs como \texttt{String}. A 100~M aristas
eso son 20~GB ---inviable en una notebook---. La
\texttt{CompactRelation} es el equivalente \emph{interno} que el
matcher itera: 32~bytes exactos, \texttt{\#[repr(C)]}, dos por
\emph{cache line}, \emph{mmap-friendly}. Este capítulo justifica
cada byte del layout, explica por qué \texttt{repr(C)} es
obligatorio (no opcional), y muestra cómo el motor convierte de
\texttt{Relation} pública a \texttt{CompactRelation} en la frontera
del ingestor.
\end{intuicion}

\section{El layout}

\begin{lstlisting}[language=Rust]
#[repr(C)]
pub struct CompactRelation {
    pub timestamp: i64,            // 8 B (offset  0)
    pub metadata_offset: u64,       // 8 B (offset  8)  handle al MetadataStore
    pub source_sid: StrId,          // 4 B (offset 16)
    pub dest_sid: StrId,            // 4 B (offset 20)
    pub ingested_at_delta: u32,     // 4 B (offset 24)  segundos desde INGEST_EPOCH_BASE
    pub dataset_tag: u16,           // 2 B (offset 28)
    pub rel_type_tag: u8,           // 1 B (offset 30)
    // 1 B padding implícito        // 1 B (offset 31)
}
// Total: 32 bytes alineados a 8.
\end{lstlisting}

\begin{figure}[H]
\centering
\begin{tikzpicture}[
    every node/.style={font=\scriptsize, align=center},
    cell/.style={draw=accent, thick, fill=accentlight, minimum height=1cm, font=\scriptsize}
]
\node[cell, minimum width=2.4cm] (ts)      {timestamp\\\scriptsize 8\,B};
\node[cell, minimum width=2.4cm, right=0pt of ts] (mo)  {metadata\_offset\\\scriptsize 8\,B};
\node[cell, minimum width=1.2cm, right=0pt of mo] (src) {src\_sid\\\scriptsize 4\,B};
\node[cell, minimum width=1.2cm, right=0pt of src] (dst) {dst\_sid\\\scriptsize 4\,B};

\node[cell, minimum width=1.2cm, below=0.5cm of ts.south west, anchor=north west]
    (ing) {ingested\_at\\\scriptsize 4\,B};
\node[cell, minimum width=0.8cm, right=0pt of ing] (ds)  {ds\_tag\\\scriptsize 2\,B};
\node[cell, minimum width=0.6cm, right=0pt of ds]  (rel) {rel\_tag\\\scriptsize 1\,B};
\node[cell, minimum width=0.6cm, right=0pt of rel, fill=white, draw=accent, dashed]
     (pad) {pad\\\scriptsize 1\,B};

\draw[decorate, decoration={brace, mirror, amplitude=4pt}, accent, thick]
    ([yshift=-2pt]ts.south west) -- ([yshift=-2pt]dst.south east)
    node[midway, below=10pt]{\footnotesize 24~B (dos cache half-lines)};
\end{tikzpicture}
\caption{\texttt{CompactRelation} en \texttt{\#[repr(C)]}: 32 bytes
exactos, alineados a 8. El \emph{layout} fue elegido para que un
\texttt{\&[StreamEdge]} (que comparte el mismo layout byte-a-byte)
se mapee directamente a una región mmap cuando el \emph{streaming
backend} \emph{spills} a disco.}
\label{fig:compact-relation-layout}
\end{figure}

\section{Por qué exactamente 32 bytes (y no 24, no 40)}

Tres restricciones cruzaron en el diseño:

\begin{enumerate}
    \item \textbf{Half cache line.} 32~B es la mitad de una cache
    line típica (64~B). Dos aristas por cache line significa que
    una iteración secuencial sobre un \texttt{Vec<CompactRelation>}
    procesa dos elementos por cada transferencia de RAM. Si pesara
    40~B, sería 1.6 elementos por línea (con líneas a medio
    aprovechar). 32~B es óptimo en \emph{packing}.
    \item \textbf{Alineación a 8 para \texttt{i64}.} El campo
    \texttt{timestamp} es \texttt{i64}; las arquitecturas exigen
    alineación a 8~B para evitar penalty (en algunas, \emph{cras h}).
    Con 32~B el siguiente \texttt{CompactRelation} arranca también
    alineado a 8 sin padding entre elementos.
    \item \textbf{Suficientes bits para los campos.} Cada campo
    tiene su requerimiento mínimo:
    \begin{itemize}
        \item \texttt{timestamp i64}: epoch en segundos cubre desde
        $\sim$292~B AC hasta $\sim$292~B DC. Margen sobrado para
        \texttt{Relation}s con timestamps razonables.
        \item \texttt{metadata\_offset u64}: handle al
        \texttt{MetadataStore}; con 64~bits podés indexar 16~exabytes
        de metadata. Margen ridículo, pero \texttt{u32} (4~B)
        limitaría a 4~GB de metadata por sesión, que en algunos
        casos extremos podría pellizcar.
        \item \texttt{source\_sid u32}, \texttt{dest\_sid u32}: el
        interner soporta hasta \(2^{32}\) strings; el motor jamás
        vio más de \(10^7\). Margen $\geq 100\times$.
        \item \texttt{ingested\_at\_delta u32}: segundos desde
        \texttt{INGEST\_EPOCH\_BASE} (epoch al inicio de la sesión).
        Cubre $\sim$136 años; el motor jamás vio una sesión más
        larga que un mes.
        \item \texttt{dataset\_tag u16}: identifica de qué dataset
        proviene la arista (sesiones multi-dataset). 65\,536
        datasets distintos por sesión es ridículamente generoso.
        \item \texttt{rel\_type\_tag u8}: 256 \emph{relation types}
        distintos. Hoy hay $\sim$15.
    \end{itemize}
\end{enumerate}

\noindent
Sumá los campos: \(8 + 8 + 4 + 4 + 4 + 2 + 1 = 31\)~B. Padding
implícito de 1~B para alinear a 8 \(\Rightarrow\) 32~B exactos.

\section{\texttt{\#[repr(C)]}: por qué es obligatorio}

Sin \texttt{\#[repr(C)]}, el compilador de Rust puede reordenar los
campos. Para un struct de 7 campos puede haber \(7! = 5040\)
permutaciones (no todas válidas por alineación, pero muchas).
Distintas versiones del compilador podrían elegir distinto. Eso
rompe tres cosas:

\begin{enumerate}
    \item \textbf{Mapeo a mmap.} El \emph{streaming backend}
    (capítulo~\ref{cap:streaming-mirror}) escribe la \emph{slice
    contiguous} \texttt{\&[StreamEdge]} (mismo \texttt{repr(C)} 32~B
    de \texttt{CompactRelation}) a un archivo durante \emph{auto-spill}
    y la mapea de vuelta a memoria. Si el layout cambia entre
    escritura y lectura ---por ejemplo, sesión guardada con
    compilador A, cargada con compilador B---, los bytes en disco
    se interpretan mal. \texttt{repr(C)} garantiza layout estable.
    \item \textbf{C ABI / FFI.} El motor expone partes de su API
    como C ABI (capítulo correspondiente del libro
    \emph{Especificación}). Pasar un \texttt{*const CompactRelation}
    a código C requiere que el layout sea predecible; sin
    \texttt{repr(C)}, no lo es.
    \item \textbf{Razonamiento sobre tamaño.} El programador eligió
    32~B; sin \texttt{repr(C)}, el compilador podría empaquetar a
    24~B o expandir a 40~B según convenga. Las cuentas de
    ``cuántas aristas caben en L2'' dejan de ser fijas.
\end{enumerate}

\begin{nota}
\texttt{repr(C)} \emph{no} garantiza endianness. Un mmap escrito en
una máquina x86\_64 little-endian y leído en una ARM big-endian
(rara hoy en día) tendría los bytes invertidos. GraphHunter asume
little-endian (todos los hosts modernos lo son); si en el futuro
hay que soportar big-endian, hay que serializar/deserializar
campo por campo en lugar de \texttt{mmap} directo.
\end{nota}

\section{Conversión: \texttt{Relation} $\to$ \texttt{CompactRelation}}

Cada \texttt{add\_relation} hace la conversión:

\begin{lstlisting}[language=Rust]
pub fn add_relation(&mut self, r: Relation) -> Result<(), Error> {
    let source_sid = self.interner.intern(&r.source_id);
    let dest_sid   = self.interner.intern(&r.dest_id);
    let metadata_offset = if r.metadata.is_empty() {
        0  // sentinel: "no metadata"
    } else {
        self.metadata_store.append(&r.metadata)?
    };
    let ingested_at_delta =
        (self.now() - self.ingest_epoch_base) as u32;
    let cr = CompactRelation {
        timestamp: r.timestamp,
        metadata_offset,
        source_sid,
        dest_sid,
        ingested_at_delta,
        dataset_tag: self.current_dataset_tag,
        rel_type_tag: r.rel_type as u8,
    };
    self.push_to_streaming(&cr);  // sole canonical store
    Ok(())
}
\end{lstlisting}

\noindent
Lo importante para el resto del libro:

\begin{itemize}
    \item \texttt{source\_sid} y \texttt{dest\_sid} salen del
    interner (capítulo~\ref{cap:interner}). Strings comunes
    (\texttt{"alice"}) se internan \emph{una vez} y todos los
    \texttt{CompactRelation} que los referencian comparten el
    \texttt{StrId}.
    \item \texttt{metadata\_offset = 0} es el sentinel
    ``sin metadata''; el matcher chequea \texttt{== 0} sin lookup
    al \texttt{MetadataStore}.
    \item \texttt{rel\_type\_tag} es un \texttt{u8} casteado del enum
    \texttt{RelationType}; el matcher usa el \texttt{u8} directo en
    sus comparaciones, sin resolver el enum.
\end{itemize}

\section{Por qué el matcher itera \texttt{CompactRelation}, no \texttt{Relation}}

El matcher hot path itera $\sim 10^7$ aristas por hunt. Para cada
arista chequea: ¿el tipo coincide?, ¿el timestamp cae en la ventana?,
¿el source binding ya está fijado? Cada chequeo tiene que ser una
operación de palabra de máquina, no una desreferencia + recorrido
de string.

\begin{itemize}
    \item Type check: \texttt{cr.rel\_type\_tag == hypothesis.step\_tag}
    (1 cache line, 1 comparación).
    \item Timestamp check: \texttt{lo $\leq$ cr.timestamp $\leq$ hi}
    (datos ya en cache desde el step anterior, 2 comparaciones).
    \item Source binding check: \texttt{cr.source\_sid ==
    binding[var]} (1 comparación de \texttt{u32}).
\end{itemize}

\noindent
Tres operaciones, tres palabras de máquina. Una iteración del
\emph{inner loop} del DFS es del orden de \(\sim 10\) ciclos.
Iterar la \texttt{Relation} pública sería del orden de \(\sim
10^3\) ciclos por la indirección a strings y al hashmap de
metadata.

\section{Resumen}

\begin{itemize}
    \item \texttt{CompactRelation}: 32~B exactos,
    \texttt{\#[repr(C)]}, alineados a 8, dos por cache line.
    \item Cada campo justificado: \texttt{i64} timestamp para
    margen, \texttt{u32} sids por el interner, \texttt{u8 rel\_tag}
    para matching sin resolver enum, \texttt{u64} metadata\_offset
    por escala extrema.
    \item \texttt{repr(C)} obligatorio para mmap stability, FFI, y
    razonamiento de tamaño.
    \item Conversión \texttt{Relation $\to$ CompactRelation} en
    \texttt{add\_relation}, vía interner para los sids y
    \texttt{MetadataStore} para metadata.
    \item El matcher itera \texttt{CompactRelation} (\(\sim 10\)
    ciclos por arista) en lugar de \texttt{Relation} (\(\sim 10^3\)
    ciclos por arista).
\end{itemize}

\begin{ejercicios}
\begin{enumerate}
    \item \starred{} Calculá: en una cache L2 de 1~MB, ¿cuántas
    aristas \texttt{CompactRelation} caben? ¿Cuántas
    \texttt{Relation} de \(200\)~B caben? ¿Cuál es el ratio?

    \item Diseñá un layout alternativo de 24~B para una
    \texttt{CompactRelation} más chica. ¿Qué campos sacrificás?
    ¿Bajo qué condiciones (qué tipo de hunts) sería preferible?

    \item \starred{} Si quitás \texttt{\#[repr(C)]}, el
    compilador podría reordenar a \texttt{(timestamp, metadata\_offset,
    source\_sid, dest\_sid, ingested\_at\_delta, dataset\_tag,
    rel\_type\_tag)} = 31~B + 1~B padding. Mismo tamaño total.
    ¿Está mal? Discutí: ¿qué se rompe específicamente?

    \item Calculá: para un grafo de \(10^8\) aristas almacenadas como
    \texttt{Vec<CompactRelation>} ordenadas, ¿cuánta memoria total
    ocupan? Si la notebook tiene 16~GB de RAM, ¿cabe?

    \item ¿Por qué \texttt{metadata\_offset} es \texttt{u64} y
    \texttt{source\_sid} es \texttt{u32}? Justificá la asimetría
    desde la escala esperada de cada uno.

    \item Imaginá que quisieras agregar un campo
    \texttt{flags: u8} a \texttt{CompactRelation} (por ejemplo,
    para marcar la arista como ``deprecated'' por una operación
    futura). ¿Cuál es el impacto en el layout? ¿Tenés que crecer
    a 40~B?
\end{enumerate}
\end{ejercicios}
